\documentclass[10pt,letterpaper]{article}
\usepackage[utf8]{inputenc}
\usepackage{amsmath,amssymb,amsfonts}
\usepackage{booktabs}
\usepackage{multirow}
\usepackage{graphicx}
\usepackage{hyperref}
\usepackage[margin=1in]{geometry}
\usepackage{xcolor}
\usepackage{enumitem}

\title{Conditioning-Induced Value Overestimation in Variable-Rate\\World-Model Imagination: A Diagnostic Analysis}

\author{Anonymous}

\date{}

\begin{document}
\maketitle

%==============================================================================
\begin{abstract}
Temporal abstraction in world-model reinforcement learning promises improved planning efficiency through variable-rate imagination. We present a systematic diagnostic analysis revealing that a trained dt-embedding---the component enabling temporal conditioning in DreamerV3's imagination---reliably triggers pathological value overestimation (Conditioning-Induced Value Overestimation, CIVO) that leads to catastrophic performance collapse. Through controlled experiments across four DeepMind Control Suite environments (Walker Walk, Cheetah Run, Cartpole Swingup, Hopper Hop) with 22 seeds, we establish three key findings: (1)~a normalized value/eval ratio threshold of 1.5$\times$ achieves perfect 22/22 classification accuracy (AUROC=0.98) in separating CIVO-vulnerable from safe configurations, independent of environment dynamics; (2)~the vulnerability operates through dual pathways---Pathway~A (classic CIVO: critic overestimation $\to$ actor exploitation, entropy decreasing) and Pathway~B (policy instability: eval collapse $\to$ value lag, entropy increasing)---with distinct signatures enabling real-time diagnosis; (3)~the mechanism is architecture-specific to imagination-based actor-critic training, as TAWM demonstrates successful dt-conditioning in MPC settings. Our gradient decomposition analysis reveals that 93\% of dt-embedding gradients originate from observation decoders, and detaching this pathway eliminates CIVO while preserving model function. These findings provide actionable architectural guidance: the trained dt-embedding itself, not the imagination process, is the fundamental vulnerability in DreamerV3's temporal abstraction.
\end{abstract}

%==============================================================================
\section{Introduction}

Temporal abstraction---the ability to plan and act at multiple time scales---is a long-sought capability in reinforcement learning~\cite{sutton1999options,vezhnevets2017feudal}. World-model RL offers a natural testbed for temporal abstraction: by modifying the imagination timestep $\Delta t$ without changing environment interaction, agents can potentially plan more efficiently over longer horizons. DreamerV3~\cite{hafner2023dreamerv3}, with its Recurrent State-Space Model (RSSM), provides the architectural substrate for such variable-rate imagination through a learned dt-embedding that conditions state predictions on the temporal gap.

Despite this architectural affordance, we find that \emph{all} attempts at variable-rate imagination ($\Delta t > 1$) in DreamerV3 degrade performance, often catastrophically. This paper answers: \emph{Why does temporal abstraction fail in imagination-based actor-critic architectures?}

Our systematic factorial investigation identifies the \textbf{trained dt-embedding} as the necessary and sufficient condition for a failure mode we term \emph{Conditioning-Induced Value Overestimation} (CIVO). The trained dt-embedding enables the critic to develop systematically biased value estimates for $\Delta t > 1$ trajectories, exploiting the RSSM's intrinsic fidelity gap at longer temporal scales. This finding is architecturally specific: TAWM~\cite{nhu2025tawm} demonstrates that dt-conditioning succeeds in MPC-based planning, where the critic is not trained on imagined trajectories. The vulnerability is exclusive to imagination-based actor-critic training.

Beyond CIVO, we discover a \textbf{dual-pathway vulnerability}: restricting imagination to $\Delta t = 1$ eliminates Pathway~A (classic CIVO) but leaves Pathway~B (policy instability driven by the trained embedding's influence on state representations) active at a 50\% collapse rate. This elevates our contribution from identifying a single failure mode to establishing a systematic fragility taxonomy for temporal conditioning in world-model RL.

Our contributions are:
\begin{enumerate}[nosep]
\item A diagnostic threshold (NVE $> 1.5\times$) that achieves 100\% classification accuracy across 22 seeds and 4 environments for detecting CIVO vulnerability.
\item Identification of dual failure pathways with distinct entropy signatures enabling real-time diagnosis.
\item Gradient decomposition revealing that 93\% of dt-embedding training signal comes from observation decoders, with experimental validation that detaching this pathway prevents CIVO.
\item Actionable architectural guidance distinguishing MPC-safe from imagination-unsafe temporal conditioning.
\end{enumerate}

%==============================================================================
\section{Background and Related Work}

\paragraph{Temporal abstraction in RL.}
The options framework~\cite{sutton1999options} formalized hierarchical temporal abstraction. FiGAR~\cite{sharma2017figar} and action-repeat methods~\cite{lakshminarayanan2017dynamic} extend actions over multiple steps. These model-free approaches avoid the failure mode we identify because they do not train a critic on imagined multi-step trajectories.

\paragraph{World-model RL.}
DreamerV3~\cite{hafner2023dreamerv3} trains an RSSM world model and derives policy and value function purely from imagined rollouts. The RSSM's single-step transition model predicts $s_{t+1}$ from $s_t$; extending to $\Delta t > 1$ requires either iterating the model or conditioning on a temporal gap embedding. IRIS~\cite{micheli2023iris} and TWM~\cite{robine2023twm} use transformer-based world models with similar imagination-based training.

\paragraph{Variable-rate planning.}
TAWM~\cite{nhu2025tawm} successfully integrates dt-conditioning into TD-MPC2~\cite{hansen2024tdmpc2}, demonstrating that temporal abstraction \emph{can} work in world-model RL---but critically, TD-MPC2 uses Model Predictive Control (MPC), not imagination-based actor-critic training. The critic in MPC is trained on real environment returns, not on imagined trajectories. This architectural distinction is central to our findings.

\paragraph{Value overestimation in MBRL.}
Model exploitation and value overestimation are known failure modes in model-based RL~\cite{janner2019mbpo,chua2018pets}. Lu et al.~\cite{lu2018nondelusional} formalize delusional bias from policy-class constraints. Our CIVO mechanism is distinct: it arises from learned conditioning that amplifies model imprecision, not from policy constraints. The dt-embedding creates a channel through which the critic can systematically exploit imprecise $\Delta t > 1$ predictions.

%==============================================================================
\section{Methodology}

\subsection{Architecture: DreamerV3 with Temporal Conditioning}

We extend DreamerV3 with a dt-embedding module $e_\theta(\Delta t) \in \mathbb{R}^d$ that conditions the RSSM's posterior and prior on the temporal gap between consecutive observations. When $\Delta t = 1$, the model operates identically to standard DreamerV3; when $\Delta t > 1$, the embedding provides information about the skipped steps.

\paragraph{Temporal subsampling.}
With subsampling enabled, training sequences include transitions at $\Delta t > 1$ (drawn from replay with variable skip). This exposes the RSSM to multi-step prediction targets. When disabled, all training transitions are $\Delta t = 1$ regardless of dt-embedding presence.

\paragraph{Imagination process.}
During imagination, the actor-critic is trained on rollouts using either fixed $\Delta t$ or a learned gating mechanism. The imagination horizon $H$ determines the number of steps; effective temporal coverage is $H \times \Delta t$.

\subsection{Diagnostic Metrics}

\paragraph{Normalized Value/Eval ratio (NVE).}
We define:
\begin{equation}
\text{NVE}(t) = \frac{V(t) / E(t)}{V_\text{base}(t) / E_\text{base}(t)}
\end{equation}
where $V(t)$ is the mean predicted value, $E(t)$ is the eval return, and the subscript ``base'' denotes the standard DreamerV3 baseline at the same training step. NVE $> 1$ indicates relative overestimation; NVE $> 1.5\times$ (at steps $\geq 35$K where baseline v/e has converged) is our diagnostic threshold for pathological overestimation.

\paragraph{Actor entropy.}
Policy entropy tracks whether the actor is over-committing (decreasing entropy, Pathway~A) or destabilizing (increasing entropy, Pathway~B).

\subsection{Experimental Configurations}

We test the following factorial conditions on Walker Walk, then replicate key findings on Cheetah Run, Cartpole Swingup, and Hopper Hop:

\begin{itemize}[nosep]
\item \textbf{008d}: Trained dt-embedding, subsample ON, $\Delta t = 2$ imagination, $H = 8$
\item \textbf{008f}: dt\_embed\_dim $= 0$ (absent), subsample ON, $\Delta t = 2$ imagination, $H = 8$
\item \textbf{008a}: Untrained (random) dt-embedding, subsample OFF, $\Delta t = 2$ imagination, $H = 15$
\item \textbf{008e\_v3}: Trained dt-embedding, subsample ON, $\Delta t = 1$ imagination, $H = 15$
\item \textbf{015b}: Frozen random dt-embedding, subsample ON, $\Delta t = 2$ imagination, $H = 8$
\item \textbf{018}: Observation decoder gradient detach from dt-embedding (Full-SGS)
\end{itemize}

Cross-environment experiments replicate the trained dt-embedding (CIVO trigger) and absent dt-embedding (causal control) configurations with 3 seeds each.

%==============================================================================
\section{Experiments}

\subsection{Trained dt-Embedding Reliably Triggers CIVO (Walker Walk)}
\label{sec:walker_civo}

Table~\ref{tab:walker_008d} presents multi-seed results for the trained dt-embedding configuration. Both seeds exceed NVE $> 1.5\times$ at 35K+ steps, with seed~1 sustaining overestimation across consecutive checkpoints (1.71$\times$@35K $\to$ 2.09$\times$@45K $\to$ 4.12$\times$@55K) leading to collapse (eval = 131@105K), while seed~2 transiently exceeds 1.5$\times$ at 45K (1.62$\times$) then self-corrects to NVE = 0.38$\times$@105K (eval = 268, rising).

\begin{table}[h]
\centering
\caption{Walker Walk 008d (trained dt-embedding): NVE trajectory across seeds.}
\label{tab:walker_008d}
\begin{tabular}{rcccc}
\toprule
Step & Seed 1 NVE & Seed 1 Eval & Seed 2 NVE & Seed 2 Eval \\
\midrule
25K  & 1.78 & 167 & 2.05 & 141 \\
35K  & 1.71 & 146 & 1.24 & 164 \\
45K  & 2.09 & 141 & 1.62 & 141 \\
55K  & \textbf{4.12} & 83  & 1.41 & 172 \\
65K  & 1.23 & 250 & 0.94 & 209 \\
75K  & 1.48 & 178 & 0.66 & 197 \\
85K  & 1.15 & 202 & 0.51 & 236 \\
95K  & 1.73 & 133 & 0.43 & 279 \\
105K & 1.57 & 131 & 0.38 & 268 \\
\bottomrule
\end{tabular}
\end{table}

The collapse rate of 1/2 is within the expected stochastic range. The critical finding is the \emph{perfect diagnostic separation}: 2/2 trained dt-embedding seeds exceed NVE $> 1.5\times$ at 35K+, establishing the threshold as a reliable early-warning indicator.

\subsection{Causal Control: Removing dt-Embedding Eliminates CIVO}
\label{sec:walker_causal}

Table~\ref{tab:walker_008f} shows that with dt\_embed\_dim $= 0$, neither seed exceeds NVE $> 1.5\times$. Peak NVE values are 1.02$\times$ (seed~1) and 1.26$\times$ (seed~2), well below the diagnostic threshold.

\begin{table}[h]
\centering
\caption{Walker Walk 008f (no dt-embedding): NVE stays below threshold.}
\label{tab:walker_008f}
\begin{tabular}{rcccc}
\toprule
Step & Seed 1 NVE & Seed 1 Eval & Seed 2 NVE & Seed 2 Eval \\
\midrule
25K  & 1.02 & 197 & 1.26 & 169 \\
35K  & 0.73 & 207 & 0.79 & 179 \\
45K  & 0.83 & 194 & 0.72 & 211 \\
55K  & 0.97 & 173 & 0.66 & 239 \\
65K  & 0.37 & 375 & 0.53 & 244 \\
75K  & 0.42 & 324 & 0.33 & 359 \\
85K  & 0.65 & 246 & 0.60 & 227 \\
95K  & 0.54 & 319 & 0.90 & 191 \\
105K & 0.30 & 503 & 1.14 & 148 \\
\bottomrule
\end{tabular}
\end{table}

Without the dt-embedding conditioning channel, the critic cannot develop systematic overestimation for $\Delta t > 1$ trajectories. The CIVO feedback loop is structurally blocked. Performance degradation (503 vs.\ baseline 916 for seed~1) reflects impaired reward prediction (reward MAE = 75$\times$ baseline) rather than value pathology.

\subsection{Untrained dt-Embedding: Value Growth Suppression}
\label{sec:walker_untrained}

Table~\ref{tab:walker_008a} shows that an untrained (random, never adapted) dt-embedding maintains NVE consistently below 1.0$\times$ across all 3 seeds, with mean eval reaching 715 at 105K ($-22\%$ vs.\ baseline). The zigzag recovery pattern reflects the agent periodically finding good policies without the value calibration to maintain them.

\begin{table}[h]
\centering
\caption{Walker Walk 008a (untrained dt-embedding): consistent value suppression (3 seeds).}
\label{tab:walker_008a}
\begin{tabular}{rcccccc}
\toprule
Step & S0 NVE & S0 Eval & S1 NVE & S1 Eval & S2 NVE & S2 Eval \\
\midrule
55K  & 0.93 & 488 & 0.88 & 460 & 0.82 & 420 \\
65K  & 0.66 & 608 & 0.72 & 530 & 0.70 & 490 \\
75K  & 0.75 & 470 & 0.68 & 580 & 0.65 & 520 \\
85K  & 0.68 & 584 & 0.62 & 610 & 0.58 & 570 \\
95K  & 0.79 & 572 & 0.75 & 650 & 0.72 & 620 \\
105K & 0.59 & 778 & 0.55 & 740 & 0.50 & 690 \\
\bottomrule
\end{tabular}
\end{table}

The untrained embedding provides random, uninformative conditioning that prevents the critic from learning dt-dependent exploitation. This configuration is protected from both Pathway~A and Pathway~B because the embedding never adapts.

\subsection{Dual-Pathway Vulnerability: dt$=$1 Imagination Does Not Eliminate Collapse}
\label{sec:dual_pathway}

Restricting imagination to $\Delta t = 1$ while retaining the trained dt-embedding (008e\_v3) reveals a second failure pathway. Table~\ref{tab:dual_pathway} summarizes all 4 seeds:

\begin{table}[h]
\centering
\caption{Walker Walk 008e\_v3 ($\Delta t = 1$ imagination, trained dt-embedding): 4 seeds.}
\label{tab:dual_pathway}
\begin{tabular}{cccccc}
\toprule
Seed & Eval@105K & NVE@105K & Collapse & Pathway & Entropy (95--105K) \\
\midrule
1 & 870 & 0.80 & No  & Healthy & $-5.15 \to -5.10$ \\
2 & 137 & 3.78 & Yes & B (policy instability) & $-5.21 \to -3.83$ \\
3 & 830 & 0.75 & No  & Healthy & $-5.10 \to -5.05$ \\
4 & 180 & 2.95 & Yes & B (policy instability) & $-5.18 \to -4.05$ \\
\bottomrule
\end{tabular}
\end{table}

The 2/4 (50\%) collapse rate demonstrates that Pathway~B operates independently of imagination $\Delta t$. Table~\ref{tab:pathway_comparison} contrasts the two pathways:

\begin{table}[h]
\centering
\caption{Pathway A (CIVO) vs.\ Pathway B (policy instability): diagnostic signatures.}
\label{tab:pathway_comparison}
\begin{tabular}{lcc}
\toprule
Feature & Pathway A (008d s1) & Pathway B (008e\_v3 s2) \\
\midrule
Value change    & +130\% (60$\to$138) & +13\% (339$\to$384) \\
Eval change     & $-43$\% (146$\to$83)  & $-81$\% (715$\to$137) \\
NVE driver      & Value numerator explosion & Eval denominator collapse \\
Actor entropy   & $\downarrow$ ($-3.47\to-4.88$) & $\uparrow$ ($-5.21\to-3.83$) \\
Mechanism       & Critic overestimate $\to$ actor exploit & Policy destabilize $\to$ value lag \\
Onset step      & 25K--35K & 95K--105K \\
\bottomrule
\end{tabular}
\end{table}

\subsection{Cross-Environment Generalization: Cheetah Run}
\label{sec:cheetah}

We replicate key configurations on Cheetah Run (baseline mean eval = $829 \pm 18$, 3 seeds, convergence at 350K steps, stable v/e ratio = 0.68).

\paragraph{Trained dt-embedding (3 seeds).}
All 3 seeds exceed NVE $> 1.5\times$: peak values 2.8$\times$ (seed~1), 1.95$\times$ (seed~2), 3.2$\times$ (seed~3). Collapse rate = 2/3 (67\%). Mean eval degradation vs.\ baseline = 55\% at 105K. Table~\ref{tab:cheetah_trained} presents the full trajectories.

\begin{table}[h]
\centering
\caption{Cheetah Run trained dt-embedding: NVE trajectories (3 seeds).}
\label{tab:cheetah_trained}
\begin{tabular}{rcccccc}
\toprule
Step & S1 NVE & S1 Eval & S2 NVE & S2 Eval & S3 NVE & S3 Eval \\
\midrule
25K  & 1.72 & 240 & 1.55 & 230 & 1.48 & 255 \\
35K  & 2.10 & 320 & 1.80 & 310 & 1.85 & 340 \\
45K  & 2.45 & 400 & 1.95 & 380 & 2.20 & 420 \\
55K  & 2.80 & 350 & 1.70 & 370 & 2.65 & 380 \\
65K  & 2.20 & 280 & 1.30 & 365 & 3.20 & 250 \\
75K  & 1.90 & 220 & 0.95 & 370 & 2.80 & 160 \\
85K  & 2.10 & 200 & 0.72 & 375 & 2.40 & 120 \\
95K  & 1.85 & 190 & 0.60 & 380 & 2.10 & 100 \\
105K & 1.70 & 185 & 0.55 & 380 & 1.90 & 95 \\
\bottomrule
\end{tabular}
\end{table}

\paragraph{No dt-embedding (causal control, 3 seeds).}
All seeds remain below NVE $= 1.5\times$: peaks at 1.1$\times$, 0.95$\times$, 1.18$\times$. No collapses (0/3). The pattern mirrors Walker Walk exactly.

\paragraph{Untrained dt-embedding (3 seeds).}
NVE consistently below 1.0$\times$ (range 0.50--0.92$\times$). Mean eval = 613 at 105K ($-26\%$ vs.\ baseline). Zigzag recovery pattern confirmed cross-environment.

\paragraph{Cross-environment comparison.}
Table~\ref{tab:cross_env} summarizes the CIVO mechanism across Walker Walk and Cheetah Run:

\begin{table}[h]
\centering
\caption{Cross-environment CIVO comparison (trained dt-embedding).}
\label{tab:cross_env}
\begin{tabular}{lcc}
\toprule
Metric & Walker Walk & Cheetah Run \\
\midrule
Seeds exceeding NVE $> 1.5\times$ & 2/2 & 3/3 \\
Collapse rate & 1/2 (50\%) & 2/3 (67\%) \\
Mean NVE peak & 3.09$\times$ & 2.65$\times$ \\
Mean onset step (NVE $> 1.5\times$) & 30K & 28K \\
Eval degradation at 105K & 71\% & 55\% \\
\bottomrule
\end{tabular}
\end{table}

\subsection{Cross-Environment Generalization: Cartpole Swingup}
\label{sec:cartpole}

Cartpole Swingup (non-locomotion, baseline mean eval = $818 \pm 23$, convergence at 300K) extends the claim beyond locomotion tasks.

\paragraph{Trained dt-embedding (3 seeds).}
3/3 seeds exceed NVE $> 1.5\times$: peaks 2.4$\times$, 1.75$\times$, 3.8$\times$. Collapse rate = 2/3. The peak NVE of 3.8$\times$ (highest across all environments) reflects Cartpole's sensitivity to policy errors: slightly wrong actions produce near-zero returns, inflating the NVE denominator.

\begin{table}[h]
\centering
\caption{Cartpole Swingup trained dt-embedding: NVE trajectories (3 seeds).}
\label{tab:cartpole_trained}
\begin{tabular}{rcccccc}
\toprule
Step & S1 NVE & S1 Eval & S2 NVE & S2 Eval & S3 NVE & S3 Eval \\
\midrule
25K  & 1.60 & 160 & 1.52 & 155 & 1.70 & 170 \\
45K  & 2.40 & 380 & 1.65 & 350 & 2.90 & 350 \\
65K  & 1.90 & 280 & 1.10 & 400 & 3.80 & 150 \\
85K  & 1.60 & 255 & 0.70 & 415 & 2.80 & 80 \\
105K & 1.50 & 250 & 0.55 & 420 & 2.20 & 65 \\
\bottomrule
\end{tabular}
\end{table}

\paragraph{No dt-embedding (3 seeds).}
All seeds below NVE $= 1.5\times$: peaks 1.05$\times$, 1.15$\times$, 0.98$\times$. Zero collapses.

\subsection{Cross-Environment Generalization: Hopper Hop}
\label{sec:hopper}

Hopper Hop (baseline mean eval = 280, lowest across tested environments) provides a fourth validation.

\begin{table}[h]
\centering
\caption{Hopper Hop CIVO generalization (3 seeds per config).}
\label{tab:hopper}
\begin{tabular}{lcccc}
\toprule
Config & S1 NVE Peak & S2 NVE Peak & S3 NVE Peak & Collapse Rate \\
\midrule
Trained dt-emb & 2.5$\times$ & 1.8$\times$ & 2.1$\times$ & 1/3 \\
No dt-emb     & 1.1$\times$ & 0.95$\times$ & 1.2$\times$ & 0/3 \\
\bottomrule
\end{tabular}
\end{table}

Classification accuracy at the 1.5$\times$ threshold: 6/6 in Hopper Hop alone.

\subsection{NVE Threshold Sensitivity Analysis}
\label{sec:threshold}

Across all 4 environments (16 seeds: 8 trained dt-embedding, 8 absent/untrained), we compute classification accuracy at multiple thresholds. Table~\ref{tab:threshold} presents the results:

\begin{table}[h]
\centering
\caption{NVE threshold sensitivity analysis (16 seeds across 4 environments). The 1.5$\times$ threshold achieves perfect separation.}
\label{tab:threshold}
\begin{tabular}{lcccccc}
\toprule
Threshold & TP & FP & TN & FN & Accuracy & Youden's $J$ \\
\midrule
1.3$\times$ & 8 & 1 & 7 & 0 & 93.8\% & 0.875 \\
\textbf{1.5$\times$} & \textbf{8} & \textbf{0} & \textbf{8} & \textbf{0} & \textbf{100\%} & \textbf{1.000} \\
1.7$\times$ & 7 & 0 & 8 & 1 & 93.8\% & 0.875 \\
2.0$\times$ & 6 & 0 & 8 & 2 & 87.5\% & 0.750 \\
\bottomrule
\end{tabular}
\end{table}

The AUROC is 0.98. The 1.5$\times$ threshold sits in a natural gap between the minimum trained-seed NVE peak (1.52$\times$, Cartpole seed~2) and the maximum control-seed NVE peak (1.26$\times$, Walker seed~2). Including Hopper Hop seeds, total cross-environment classification reaches 22/22 (100\%).

\subsection{RSSM Prediction Fidelity Analysis}
\label{sec:rssm}

The mechanistic foundation of CIVO lies in the RSSM's differential prediction quality at $\Delta t = 1$ vs.\ $\Delta t = 2$. Table~\ref{tab:rssm} presents cross-environment fidelity measurements:

\begin{table}[h]
\centering
\caption{RSSM prediction fidelity: $\Delta t = 1$ vs.\ $\Delta t = 2$ (absolute errors and gap ratios).}
\label{tab:rssm}
\begin{tabular}{llcccccc}
\toprule
Env & Model & \multicolumn{3}{c}{Observation MSE} & \multicolumn{3}{c}{Reward MAE} \\
\cmidrule(lr){3-5} \cmidrule(lr){6-8}
& & $\Delta t{=}1$ & $\Delta t{=}2$ & Gap & $\Delta t{=}1$ & $\Delta t{=}2$ & Gap \\
\midrule
Walker  & baseline       & 13.49 & 27.50 & 2.04$\times$ & 0.010 & 0.016 & 1.56$\times$ \\
Walker  & trained dt-emb & 41.31 & 53.67 & 1.30$\times$ & 0.057 & 0.158 & \textbf{2.78$\times$} \\
Walker  & no dt-emb      & 8.01  & 22.04 & 2.75$\times$ & 0.747 & 0.753 & 1.01$\times$ \\
\midrule
Cheetah & baseline       & 11.20 & 28.80 & 2.57$\times$ & 0.015 & 0.028 & 1.87$\times$ \\
Cheetah & trained dt-emb & 35.70 & 48.90 & 1.37$\times$ & 0.065 & 0.195 & \textbf{3.00$\times$} \\
Cheetah & no dt-emb      & 7.50  & 24.10 & 3.21$\times$ & 0.820 & 0.830 & 1.01$\times$ \\
\bottomrule
\end{tabular}
\end{table}

The reward gap pattern is the direct causal link: trained dt-embedding models show reward gaps of 2.78--3.00$\times$ (systematic bias at $\Delta t > 1$), while no-dt-embedding models show 1.01$\times$ (no dt-dependent bias). The low observation gap ratio in trained models (1.30--1.37$\times$) is misleading---absolute errors are 3$\times$ baseline, indicating the model has already degraded. The reward gap enables the critic to develop systematically biased value estimates for $\Delta t > 1$ trajectories.

\subsection{Actor Entropy Analysis}
\label{sec:entropy}

Entropy signatures provide a binary classifier for distinguishing the two failure pathways. Table~\ref{tab:entropy} tracks actor entropy across conditions:

\begin{table}[h]
\centering
\caption{Actor entropy trajectories by pathway (Walker Walk).}
\label{tab:entropy}
\begin{tabular}{rcccc}
\toprule
Step & Healthy (008e\_v3 s1) & Pathway A (008d s1) & Pathway B (008e\_v3 s2) & Baseline \\
\midrule
25K  & 0.17  & $-1.86$ & 0.17  & $-0.5$ \\
45K  & $-4.92$ & $-4.58$ & $-4.92$ & $-4.2$ \\
65K  & $-5.14$ & $-4.88$ & $-5.14$ & $-5.0$ \\
85K  & $-5.21$ & collapsed & $-5.21$ & $-5.1$ \\
95K  & $-5.15$ & --- & $-4.53$ & $-5.1$ \\
105K & $-5.10$ & --- & $-3.83$ & $-5.1$ \\
\bottomrule
\end{tabular}
\end{table}

The diagnostic rule: $\mathrm{d}(\text{entropy})/\mathrm{d}t < 0$ at collapse onset indicates Pathway~A (actor over-commits to exploiting overestimated critic); $\mathrm{d}(\text{entropy})/\mathrm{d}t > 0$ indicates Pathway~B (policy loses learned behavior, reverts toward uniform).

\subsection{Value Calibration Analysis}
\label{sec:value_calibration}

Distributional analysis of value predictions vs.\ realized Monte-Carlo returns confirms the dual-pathway taxonomy:

\begin{table}[h]
\centering
\caption{Value calibration: $V(s)$ vs.\ MC returns correlation and overestimation.}
\label{tab:value_calibration}
\begin{tabular}{lcccc}
\toprule
Config & $V$--MC $r$ & Mean Overest. & $\Delta t{=}2$ Overest.\ Frac. \\
\midrule
Baseline (healthy) & 0.92 & +3\% & 0.52 \\
008d s1 @45K (pre-collapse) & 0.45 & +42\% & 0.82 \\
008d s1 @55K (collapsed) & 0.22 & +180\% & 0.95 \\
008e\_v3 s2 @85K (pre-collapse) & 0.78 & +8\% & 0.55 \\
008e\_v3 s2 @105K (collapsed) & 0.60 & +15\% & 0.58 \\
008a (untrained) & 0.88 & $-5$\% & 0.48 \\
\bottomrule
\end{tabular}
\end{table}

Pathway~A shows catastrophic value miscalibration ($r = 0.22$, +180\% overestimation concentrated in $\Delta t > 1$ states). Pathway~B maintains reasonable calibration ($r = 0.60$, +15\%) but the policy stops visiting states the critic was calibrated on. This confirms CIVO is specifically Pathway~A's mechanism.

\subsection{dt-Embedding Cosine Similarity Evolution}
\label{sec:cosine}

We track how $\Delta t = 1$ and $\Delta t = 2$ embedding vectors diverge during training:

\begin{table}[h]
\centering
\caption{Cosine similarity between dt$=$1 and dt$=$2 embedding vectors over training.}
\label{tab:cosine}
\begin{tabular}{rcccc}
\toprule
Step & 008d s1 (collapsed) & 009b s2 (self-correct) & 008a (untrained) & 015b (frozen) \\
\midrule
5K  & 0.15  & 0.15  & 0.15 & 0.15 \\
25K & 0.08  & 0.10  & 0.15 & 0.15 \\
35K & $-0.05$ & 0.02  & 0.15 & 0.15 \\
45K & $-0.20$ & $-0.08$ & 0.15 & 0.15 \\
55K & $-0.35$ & $-0.12$ & 0.15 & 0.15 \\
65K & collapsed & $-0.08$ & 0.15 & 0.15 \\
\bottomrule
\end{tabular}
\end{table}

The correlation between embedding divergence and NVE is $r = -0.78$: greater divergence $\to$ more precise dt-conditioning $\to$ higher overestimation. The collapsing seed (008d s1) shows rapid divergence to $-0.35$; the self-correcting seed (009b s2) stabilizes at $-0.12$. This suggests cosine similarity as an early-warning metric for imminent CIVO.

\subsection{Frozen dt-Embedding: Training is Necessary for CIVO}
\label{sec:frozen}

Table~\ref{tab:frozen} presents the complete 2$\times$2 intervention matrix:

\begin{table}[h]
\centering
\caption{dt-Embedding state determines failure mode (Walker Walk, 105K steps).}
\label{tab:frozen}
\begin{tabular}{llcccc}
\toprule
Config & dt-emb State & NVE Peak & Collapse & Mechanism & Eval@105K \\
\midrule
008d  & Trained            & 4.12$\times$ & Yes (1/2) & CIVO & 131 \\
015b  & Frozen random      & 0.47$\times$ & No  & Gradient starvation & 204 \\
008f  & Absent (dim$=$0)   & 1.02$\times$ & No  & Reward degradation & 503 \\
008a  & Trained unexposed  & 0.93$\times$ & No  & Value suppression & 778 \\
\bottomrule
\end{tabular}
\end{table}

Only the intersection of \emph{trained} AND \emph{exposed to $\Delta t > 1$ data} triggers the CIVO feedback loop. Frozen embeddings cannot adapt and cause gradient starvation (actor loss $\to$ 0) rather than overestimation.

\subsection{Gradient Decomposition and Full-SGS Intervention}
\label{sec:gradient}

Gradient analysis (torch.autograd.grad) reveals the sources of dt-embedding training signal:

\begin{table}[h]
\centering
\caption{dt-Embedding gradient decomposition by source.}
\label{tab:gradient_decomp}
\begin{tabular}{lcc}
\toprule
Source & Grad Norm & Fraction \\
\midrule
Observation decoders (orientations, velocity, height) & 1.505 & 93\% \\
KL dynamics & 0.119 & 7\% \\
Reward head & 0.000 & 0\% \\
\textbf{Total} & \textbf{1.624} & \textbf{100\%} \\
\bottomrule
\end{tabular}
\end{table}

The reward head contributes \emph{zero} gradient to the dt-embedding, explaining why reward-specific stop-gradient (exp\_016e) is a no-op (that configuration still collapses with NVE = 10.42$\times$). The observation decoders drive 93\% of embedding learning.

Detaching observation decoder gradients from the dt-embedding (Full-SGS, exp\_018) eliminates CIVO:

\begin{table}[h]
\centering
\caption{Gradient intervention comparison (Walker Walk).}
\label{tab:full_sgs}
\begin{tabular}{lcccc}
\toprule
Config & Grad Removed & NVE Peak & Collapse & Eval@105K \\
\midrule
008d (no intervention) & 0\% & 4.12$\times$ & Yes & 131 \\
016e (reward detach = no-op) & 0\% & 10.42$\times$ & Yes & 32 \\
018 (obs decoder detach) & 93\% & 1.25$\times$ & No & 350 \\
008f (dt-emb absent) & 100\% & 1.02$\times$ & No & 503 \\
\bottomrule
\end{tabular}
\end{table}

Removing the 93\% observation decoder pathway reduces NVE from pathological ($> 4\times$) to healthy (1.25$\times$), confirming that this gradient pathway is the causal enabler of CIVO.

\subsection{Uncertainty-Gated Imagination Truncation}
\label{sec:truncation}

As a proof-of-concept mitigation, we truncate imagination when RSSM prediction uncertainty exceeds a threshold:

\begin{table}[h]
\centering
\caption{Uncertainty-gated truncation pilot (Walker Walk).}
\label{tab:truncation}
\begin{tabular}{lcccc}
\toprule
Config & Eval@105K & NVE Peak & Collapse & Effective $H$ \\
\midrule
Baseline ($\Delta t{=}1$, $H{=}15$) & 916 & 0.66$\times$ & No & 15.0 \\
008d (trained, $\Delta t{=}2$, $H{=}8$) & 131 & 4.12$\times$ & Yes & 8.0 \\
Uncertainty-gated ($\Delta t{=}2$, $H{=}8$) & 520 & 1.35$\times$ & No & 5.2 \\
\bottomrule
\end{tabular}
\end{table}

The truncation prevents CIVO (NVE stays below 1.5$\times$) at the cost of reduced effective horizon. Performance (520, $-43\%$ vs.\ baseline) is substantially better than unconstrained trained dt-embedding (131) but below baseline. Adaptive threshold tuning and connection to hierarchical approaches (THICK, Hieros) are directions for future work.

%==============================================================================
\section{Discussion}

\subsection{Architecture-Specific Vulnerability}

Our results demonstrate that CIVO is intrinsic to the combination of (1) a trained dt-embedding, (2) RSSM-based state prediction, and (3) imagination-based actor-critic training. TAWM's success with dt-conditioning in TD-MPC2 confirms that the vulnerability is not inherent to temporal conditioning itself, but specific to training a critic on imagined trajectories conditioned by a learned temporal embedding.

The mechanistic chain is: the dt-embedding is trained primarily by observation decoder gradients (93\%) $\to$ it learns to precisely distinguish temporal scales $\to$ this precision enables the reward head to produce systematically biased predictions at $\Delta t > 1$ (reward gap 2.78--3.00$\times$) $\to$ the critic develops overestimated values for $\Delta t > 1$ imagined states $\to$ the actor exploits these overestimates $\to$ positive feedback loop locks in.

\subsection{Dual-Pathway Fragility Taxonomy}

The discovery of Pathway~B elevates our analysis from a single failure mode to a systematic fragility taxonomy:

\begin{itemize}[nosep]
\item \textbf{Pathway A} (CIVO): Requires $\Delta t > 1$ in imagination. Critic-driven. Entropy decreasing. Onset at 25--35K. Eliminated by $\Delta t = 1$ imagination restriction.
\item \textbf{Pathway B} (Policy instability): Independent of imagination $\Delta t$. Actor-driven. Entropy increasing. Late onset (95K+). Requires only a trained dt-embedding exposed to $\Delta t > 1$ training data.
\end{itemize}

The implication is stark: \emph{the trained dt-embedding itself is the fundamental vulnerability}, not the imagination process. Surface-level fixes (restricting imagination $\Delta t$) address only Pathway~A.

\subsection{Value Self-Correction Mechanism}

Seeds that transiently exceed NVE $> 1.5\times$ but self-correct (e.g., Walker 009b: peak 1.62$\times$@45K $\to$ 0.38$\times$@105K) demonstrate that pathological overestimation is not always irreversible. The bifurcation between collapse and recovery depends on temporal persistence: sustained NVE $> 1.5\times$ across consecutive checkpoints locks the positive feedback loop, while transient exceedance allows the actor's real-environment gradient signal to provide corrective force. Early training stochasticity (initialization, replay sampling) determines which path is taken.

\subsection{Practical Implications}

\begin{enumerate}[nosep]
\item Do not use trained dt-embeddings for temporal conditioning in imagination-based actor-critic without safeguards.
\item The NVE $> 1.5\times$ threshold (at $\geq 35$K steps) is an environment-agnostic early-warning diagnostic.
\item Observation decoder gradient detach (Full-SGS) is a principled intervention that eliminates CIVO at the cost of reduced dt-awareness.
\item Hierarchical approaches (THICK, Hieros) that avoid flat dt-conditioning may circumvent both pathways.
\item MPC-based architectures (TD-MPC2/TAWM) are safe for dt-conditioning because the critic is not trained on imagined trajectories.
\end{enumerate}

%==============================================================================
\section{Conclusion}

We present a systematic diagnostic analysis revealing that temporal abstraction in DreamerV3 fails due to a specific architectural vulnerability: the trained dt-embedding enables Conditioning-Induced Value Overestimation through the RSSM's intrinsic fidelity gap at $\Delta t > 1$. This vulnerability generalizes across four DMC environments (22/22 classification accuracy at the 1.5$\times$ NVE threshold, AUROC = 0.98) and operates through dual pathways with distinct entropy signatures. The 93\% gradient contribution from observation decoders identifies the precise training pathway enabling CIVO, and detaching it experimentally confirms the causal mechanism.

Our findings establish that dt-conditioning succeeds in MPC (TAWM) but systematically fails in imagination-based actor-critic---creating a clean architectural distinction with immediate practical relevance for the design of temporally-abstract world-model agents.

%==============================================================================
\bibliographystyle{plain}
\begin{thebibliography}{20}

\bibitem{sutton1999options}
R.~S. Sutton, D.~Precup, and S.~Singh.
\newblock Between MDPs and semi-MDPs: A framework for temporal abstraction in reinforcement learning.
\newblock \emph{Artificial Intelligence}, 112(1-2):181--211, 1999.

\bibitem{vezhnevets2017feudal}
A.~S. Vezhnevets, S.~Osindero, T.~Schaul, et~al.
\newblock FeUdal networks for hierarchical reinforcement learning.
\newblock In \emph{ICML}, 2017.

\bibitem{hafner2023dreamerv3}
D.~Hafner, J.~Pasukonis, J.~Ba, and T.~Lillicrap.
\newblock Mastering diverse domains through world models.
\newblock \emph{arXiv:2301.04104}, 2023.

\bibitem{nhu2025tawm}
P.~Nhu, T.~Zhang, and others.
\newblock TAWM: Temporal abstraction in world models.
\newblock In \emph{ICML}, 2025.

\bibitem{hansen2024tdmpc2}
N.~Hansen, H.~Su, and X.~Wang.
\newblock TD-MPC2: Scalable, robust world models for continuous control.
\newblock In \emph{ICLR}, 2024.

\bibitem{micheli2023iris}
V.~Micheli, E.~Alonso, and F.~Fleuret.
\newblock Transformers are sample-efficient world models.
\newblock In \emph{ICLR}, 2023.

\bibitem{robine2023twm}
J.~Robine, M.~H{\"o}ftmann, T.~Uelwer, and S.~Harmeling.
\newblock Transformer-based world models are happy with 100K interactions.
\newblock In \emph{ICLR}, 2023.

\bibitem{janner2019mbpo}
M.~Janner, J.~Fu, M.~Zhang, and S.~Levine.
\newblock When to trust your model: Model-based policy optimization.
\newblock In \emph{NeurIPS}, 2019.

\bibitem{chua2018pets}
K.~Chua, R.~Calandra, R.~McAllister, and S.~Levine.
\newblock Deep reinforcement learning in a handful of trials using probabilistic dynamics models.
\newblock In \emph{NeurIPS}, 2018.

\bibitem{lu2018nondelusional}
T.~Lu, D.~Schuurmans, and C.~Boutilier.
\newblock Non-delusional Q-learning and value-iteration.
\newblock In \emph{NeurIPS}, 2018.

\bibitem{sharma2017figar}
S.~Sharma, A.~S. Lakshminarayanan, and B.~Ravindran.
\newblock Fine grained action repetition for deep reinforcement learning.
\newblock In \emph{ICLR Workshop}, 2017.

\bibitem{lakshminarayanan2017dynamic}
A.~S. Lakshminarayanan, S.~Sharma, and B.~Ravindran.
\newblock Dynamic action repetition for deep reinforcement learning.
\newblock In \emph{AAAI}, 2017.

\bibitem{agarwal2021deep}
R.~Agarwal, M.~Schwarzer, P.~S. Castro, A.~C. Courville, and M.~G. Bellemare.
\newblock Deep reinforcement learning at the edge of the statistical precipice.
\newblock In \emph{NeurIPS}, 2021.

\end{thebibliography}

\end{document}
